\section*{Erd\H{o}s Problem \#348: robustness of complete sequences}
\subsection*{1. Statement and conventions}
For which integers $0\le m<n$ does an indexed nondecreasing sequence
of positive integers remain complete after every deletion of $m$ terms,
but fail to be complete after every deletion of $n$ terms?
Repeated values have separate indices. Here \emph{complete} means that
finite subset sums contain every sufficiently large integer; strong
completeness means that they contain every nonnegative integer.
We prove the standard weak examples with an explicit infinite-hole
argument, and reconstruct van Doorn's impossibility result for the
strong interpretation. The intended weak problem for $m\ge2$ remains
unresolved.

\subsection*{2. The elementary strong-completeness criterion}
If $a_1\le a_2\le\cdots$ are positive integers and
$S_r=\sum_{i\le r}a_i$, strong completeness is equivalent to
\[
 a_1=1,\qquad a_{r+1}\le S_r+1\quad(r\ge1).
\]
Sufficiency follows by joining the intervals $[0,S_r]$ and
$[a_{r+1},a_{r+1}+S_r]$ inductively. Necessity follows because otherwise
$1$, or $S_r+1$, lies below all later terms and exceeds the sum of all
earlier terms. This argument is about indexed multisets as well.

\subsection*{3. Full verification of the weak examples}
The sequence $1,2,4,\ldots$ is strongly complete. Deleting $2^t$
excludes every integer whose binary expansion has digit $1$ in position
$t$; this is an infinite residue obstruction. Thus it gives $(m,n)=(0,1)$.

Let $F_1=F_2=1$, $F_{j+2}=F_{j+1}+F_j$. Recall
\[
 \sum_{i=1}^j F_i=F_{j+2}-1.
\]
After deleting $F_t$, the remaining sequence is strongly complete.
Indeed, before the deletion the criterion above is unchanged.
For a retained $F_j$ with $j>t$, the earlier available sum is
$F_{j+1}-1-F_t$, so the necessary inequality reduces to
$F_t\le F_{j-1}$. The initial available term is still $1$.

Deleting \emph{any two} terms destroys even weak completeness.
Fix deleted indices $a<b$ and write $D=F_a+F_b$. For each
$1\le r\le F_a$, the integer
\[
 H_b=F_{b+1}-r
\]
is missing: all available terms through index $b$ sum to
$F_{b+1}-1-F_a<H_b$, while every later term is at least $F_{b+1}$.
Suppose $H_j=F_{j+1}-r$ is missing for some $j\ge b$.
Consider $H_{j+2}=F_{j+3}-r$. No term beyond index $j+2$ can occur.
If $F_{j+2}$ is used, its removal would represent the previous hole
$F_{j+1}-r$. If it is not used, the sum of all available terms through
$j+1$ is
\[
 F_{j+3}-1-D<F_{j+3}-r.
\]
Both possibilities fail. The infinitely many values $H_{b+2t}$ are
therefore missing. This completes the weak $(1,2)$ example, including
the point that a single small hole would not have been enough.
Further deletions preserve these obstructions, giving every $(0,n)$
with $n\ge1$ and every $(1,n)$ with $n\ge2$.

\subsection*{4. Why the strong interpretation has no examples for $m\ge2$}
Suppose deletion of any two terms preserves strong completeness.
For $r\ge2$, delete the terms at indices $r-1,r$. Applying the necessary
criterion just before the retained term $a_{r+1}$ gives
\[
 a_{r+1}\le S_{r-2}+1,\qquad
 S_r+1-a_{r+1}\ge a_{r-1}+a_r.                    \tag{1}
\]
If the sequence is unbounded, choose any prescribed number $n$ of
indices $3\le i_1<\cdots<i_n$, successively far enough apart that
\[
 a_{i_j-1}\ge\sum_{\ell<j}a_{i_\ell}.
\]
After these deletions, let $D_r$ denote the sum deleted through index
$r$. If the last such index is $i_j\le r$, then
\[
 D_r\le a_{i_j-1}+a_{i_j}\le a_{r-1}+a_r.
\]
Thus (1) implies $a_{r+1}\le S_r-D_r+1$ at every retained term
$a_{r+1}$. Before the first deletion the original inequalities hold,
and the first term was not removed. The remaining sequence is strongly
complete. In particular these $n$ deletions do not destroy completeness.

If instead the nondecreasing integer sequence is bounded, it has an
infinite constant tail. Removing any fixed number of terms sufficiently
far out in that tail leaves the subset-sum set unchanged, because every
finite use of that value can be replaced by unused indices.
The same conclusion follows.

Finally, robustness under all $m\ge2$ deletions implies robustness under
all two deletions: extend any two deleted indices to $m$ indices, and
then restore the extra terms. This proves the claimed impossibility
for all $2\le m<n$ in the strong interpretation.

\subsection*{5. Assessment and references}
The weak examples and the full strong classification are proved here.
The decisive inequality (1) need not hold for a sequence with finitely
many exceptional subset sums, so this does not settle the intended
weak $(2,3)$ case or the general weak classification.

Sources: \texttt{https://www.erdosproblems.com/348} and Wouter van Doorn's
original discussion at
\texttt{https://www.erdosproblems.com/forum/thread/348}, checked 2026-09-25.
The strong argument reconstructs that known result, with the bounded
case and finite-deletion quantifiers explicit. Earlier GPT Pro 5.2
work was checked; the infinite-hole proof avoids equating failure of
strong completeness with failure of weak completeness.
No novelty or formal verification is claimed.
\subsection*{COMPLETION ESTIMATE}
\textbf{COMPLETION: 65\%}
